\section{Resumo}

% Objetivos, Metodologia, Indicação sobre os valores medidos, incluindo a sua precisão e exatidão. Identificar os objetivos atingidos.

O segundo trabalho teve como objetivo montar um circuito mais complexo, um circuito RC, e analisar como se processa a carga e a descarga de um condensador. Outro objetivo passou por determinar os intervalos de tempo associados a cada um destes processos.
Para tal procedeu-se à medição dos valores de tensão nos terminais do condensador e do intervalo de tempo para cada valor de tensão, recorrendo a um multímetro e a uma gravação de vídeo que permitia visualizar os valores de tensão mostrados no multímetro e o instante em que estes se verificavam.
\paragraph{}Com os resultados obtidos, psteriormente procedeu-se à linearização das equações teóricas do circuito RC, o que permitiu determinar as respetivas constantes de tempo. Assim, obtiveram-se os valores teóricos $\tau_c = 6\,s$ e $\tau_d = 10\,s$, bem como os valores experimentais $\tau_c = 10.30\,s$ e $\tau_d = 12.25\,s$. A comparação entre os valores teóricos e experimentais permitiu avaliar a concordância entre o modelo teórico e o comportamento observado experimentalmente.